\section{Tableau combinatorics}
We now examine some tableau combinatorics which we shall use to help us understand the double cosets of certain pairs of subgroups in
$S_n$. The material in this section is taken from the account given by Wildon in his unpublished note \cite{WLDDCR}.

Throughout this section we fix $\alpha = (\alpha_1,\alpha_2,\ldots,\alpha_l)$ to be a composition of $n$ of length $l$, and
$\gamma = (\gamma_1,\gamma_2,\ldots,\gamma_t)$ to be a composition of $n$ of length $t$.
For our fixed composition $\alpha$ of $n$, we let $S_n$ act (from the right) on both the set of tableaux of shape $\alpha$ and type
$\gamma$, by permuting the entries of a tableau as follows.
For a tableau $\tau$, number the boxes of the tableau from 1 to $n$ going from left to right across each row in turn, starting with the
top row and working down. Let $\sigma \in S_n$. Then $\tau\sigma$ is defined to be the tableau obtained from $\tau$ by moving the number
in box number $i$ to box number $(i)\sigma$, for each $i=1,\ldots,n$. For example, let us take $n=13$,
$\alpha = (5,3,4,1)$, $\gamma = (4,5,4)$, $\sigma = (1,12,3,6)(5,7,13)(8,10) \in S_{13}$, and
\[
\tau \quad=\quad 
\mytab{
{{1,2,1,3,2},
{2,3,2},
{2,3,1,3},
{1}}
}\quad.
\]
The reader may verify that we have
\[
\tau\sigma \quad=\quad 
\mytab{
{{2,2,3,3,1},
{1,2,3},
{2,2,1,1},
{3}}
}\quad.\]

It is easy to see that this definition does indeed yield a $S_n$ action as claimed, and it is obvious that this $S_n$ action is
transitive. It is natural to ask what the stabilizer of a given tableau is under this action, and in order to answer this we now consider
certain special tableaux of shape $\alpha$ and type $\gamma$. Indeed, for our compositions $\alpha$ and $\gamma$, we
construct the \textit{standard tableau of shape $\alpha$ and type $\gamma$} as follows: we begin with a Young diagram of shape $\alpha$
with the boxes numbered as described above, and then working from box 1 to box $n$ we enter first $\gamma_1$ 1's, then $\gamma_2$ 2's,
and so on. We denote this tableau by $\tau^\alpha_{\gamma}$. For example, if we take
$n=13$, $\alpha = (2,0,3,1,3,4)$ and $\gamma=(3,5,0,4,1)$, then we have
\[
\tau^\alpha_{\gamma}\quad=\quad
\mytab{
{{1,1},
{},
{1,2,2},
{2},
{2,2,4},
{4,4,4,5}}
}.
\]

\begin{prop}\label{tab_stab:prop}
For any $\sigma \in S_n$, we have $\mathrm{Stab}\bigl(\tau^\alpha_\gamma\sigma\bigr) = \bigl(S_\gamma\bigr)^\sigma$,
where we write $\mathrm{Stab}(-)$ to denote a stabilizer.
\end{prop}
\begin{proof}Elementary.\end{proof}

We say that a tableau of shape $\alpha$ and type $\gamma$ has \textit{weakly increasing rows} if the entries in its rows are weakly
increasing from left to right. Define $\mathcal{W}^\alpha_{\gamma}$ to be the set of all tableaux of shape $\alpha$ and type
$\gamma$ with weakly increasing rows.

Recall that the \textit{length} of a permutation is defined to be the total number of inversions of the permutation, where an
\textit{inversion} of a permutation $\sigma \in S_n$ is a pair $(i,j)$ such that $1 \leq i < j \leq n$ and $(i)\sigma > (j)\sigma$.
Let us take $\Omega^\alpha_{\gamma}$ to be a complete system of $(S_\gamma,S_\alpha)$-double coset representatives in $S_n$, where each element
$\sigma$ of $\Omega^\alpha_\gamma$ is of minimal length in its left coset $\sigma S_\alpha$.

\begin{prop}\label{tab_dcos_bij:prop}\textup{(}\cite{WLDDCR}, Theorem 4.1 and Corollary 5.1\textup{)}
If $\sigma \in S_n$ is of minimal length in its left $S_\alpha$-coset $\sigma S_\alpha$, then $\tau^\alpha_{\gamma}\sigma$ has
weakly increasing rows. Thus we have a map
$\Omega^\alpha_\gamma \longrightarrow \mathcal{W}^\alpha_\gamma$, $\sigma \longmapsto \tau^\alpha_\gamma\sigma$, and this map is in fact a
bijection.
\end{prop}

\begin{coro}\label{tabs_give_cosets:cor}
Suppose that we have $\sigma_1,\ldots,\sigma_N \in S_n$ such that if $i \neq j$ then
$\tau^\alpha_\gamma\sigma_i \neq \tau^\alpha_\gamma\sigma_j$ and further $\{\tau^\alpha_\gamma\sigma_i \mid 1 \leqslant i \leqslant N\} =
\mathcal{W}^\alpha_\gamma$. Then $\sigma_1,\ldots,\sigma_N$ is a complete system of $(S_\gamma,S_\alpha)$-double coset representatives in
$S_n$ without redundancy.
\end{coro}
\begin{proof}
With our system of $(S_\gamma,S_\alpha)$-double coset representatives $\Omega^\alpha_{\gamma}$ as above, we may by Proposition
\ref{tab_dcos_bij:prop} list the distinct elements of $\Omega^\alpha_{\gamma}$ as $\omega_1,\ldots,\omega_N$ such that
$\tau^\alpha_\gamma\sigma_i = \tau^\alpha_\gamma\omega_i$. This implies that $\tau^\alpha_\gamma = \tau^\alpha_\gamma\omega_i\sigma_i^{-1}$, and
hence that $\omega_i\sigma_i^{-1} \in \mathrm{Stab}(\tau^\alpha_\gamma)$, so that by Proposition \ref{tab_stab:prop} we have
$\omega_i\sigma_i^{-1} \in S_\gamma$. Hence $S_\gamma \sigma_i S_\alpha = S_\gamma (\omega_i\sigma_i^{-1}) \sigma_i S_\alpha =
S_\gamma \omega_i S_\alpha$, and so $\sigma_1,\ldots,\sigma_N$ is a complete system of $(S_\gamma,S_\alpha)$-double coset representatives in
$S_n$ without redundancy.
\end{proof}
